% ======================================================================
% SECTION 11: SITE INTERFACE LAYER AND INVESTIGATOR COORDINATION
% File: sections/site_interface.tex
% ======================================================================

[PROCESSING INSTRUCTION: Generate this section (approximately 2-3 pages) covering the
site interface layer between the autonomous sponsor and trial sites. Use the following:

SOURCE FILES TO PROCESS:
- sponsor/input\_files/sponsor\_04\_operations\_enrollment.md (site selection, contracting)
- sponsor/input\_files/org\_04\_functions\_site\_enrollment.md (RWD-driven feasibility, startup)
- sponsor/input\_files/org\_07\_cro\_split\_references.md (site-related sponsor vs. CRO split)
- national-platform/new\_paper/final\_paper/sections/site\_establishment.tex (site documentation)
- national-platform/new\_paper/final\_paper/sections/patient\_instructions.tex (patient interface)
- national-platform/new\_paper/final\_paper/sections/national\_mcp.tex (MCP site connectivity)
- national-platform/new\_trial\_psl/ (11 trial site documents)

SUBSECTION STRUCTURE:

\subsection{Site Gateway (site\_gateway)}
- Site qualification and readiness assessment
- Training module delivery and competency tracking
- Scheduling and resource allocation
- Robot-task authorization management
- Data exchange interface with site systems (EHR, PACS, LIS)

\subsection{Investigator Interface}
- Protocol training and certification automation
- Delegation of authority log management
- Medical monitoring communication channels
- Adverse event reporting interface
- Investigator oversight dashboard

\subsection{Site Systems Integration}
- EHR integration via FHIR R4 APIs
- PACS/DICOM integration for imaging
- Laboratory Information System connectivity
- Pharmacy system interface for IMP management
- MCP server endpoints for each site

\subsection{Patient-Facing Interface}
- eConsent platform management
- Patient scheduling and appointment automation
- ePRO data collection and monitoring
- Patient communication and education
- Digital health technology integration

PYTHON SCRIPT REQUIREMENTS (to be generated by Claude Code Opus 4.6 in the future):
- File: sponsor/template/scripts/site\_gateway.py
  Purpose: Site interface agent. Implement site qualification scoring
  (using PSL framework from national platform), training module delivery,
  scheduling optimization, robot-task authorization, and data exchange
  with site EHR/PACS/LIS systems via FHIR/DICOM protocols.
- File: sponsor/template/scripts/investigator\_interface.py
  Purpose: Investigator-facing interface. Protocol training automation,
  delegation log management, AE reporting workflows, and medical
  monitoring communication. Include competency assessment tracking.
- File: sponsor/template/scripts/site\_systems\_connector.py
  Purpose: Multi-protocol site systems connector. Implement FHIR R4 client
  for EHR queries, DICOM C-FIND/C-MOVE for imaging, HL7 v2 for lab results,
  and MCP tool endpoints for sponsor-site data exchange.
- File: sponsor/template/scripts/patient\_interface.py
  Purpose: Patient-facing interface manager. eConsent delivery and tracking,
  appointment scheduling, ePRO questionnaire deployment, patient education
  content delivery, and DHT data collection configuration.
- All scripts must pass ruff lint (line-length 120, E/F/W rules).
- Process with Claude Code Opus 4.6 (1M token context).

TABLE REQUIREMENTS:
- Table 19: Site Systems Integration Matrix
  Columns: System | Protocol | Data Type | Direction | Agent Handler | Security
- Table 20: Site Readiness Assessment Criteria (PSL-based)
  Columns: Category | Criterion | Scoring Method | Minimum Threshold | Evidence Required

FORMATTING NOTES:
- Reference the 11 trial site documents from national-platform/new\_trial\_psl/.
- Cross-reference \S\ref{sec:robotics} for robot-task authorization.
- Cross-reference \S\ref{sec:data} for data exchange protocols.]

\subsection{Site Gateway}

[PROCESSING INSTRUCTION: Generate 1 page on the site\_gateway. This agent manages all
interactions between the autonomous sponsor and trial sites. Cover: site qualification
using PSL framework (3 dimensions from the national platform), training module delivery
with competency verification, scheduling optimization across multiple sites, robot-task
authorization (the sponsor must authorize each robotic procedure at a site), and data
exchange via FHIR/DICOM/MCP. Reference the 11 trial site documents from
national-platform/new\_trial\_psl/ for site requirements. Reference the site establishment
section of the national platform paper. Process with Claude Code Opus 4.6
(1M token context).]

\subsection{Investigator Interface}

[PROCESSING INSTRUCTION: Generate 0.75 page on the investigator interface. Cover protocol
training automation, delegation of authority management, medical monitoring channels,
AE reporting workflows, and the investigator oversight dashboard. Reference ICH E6(R3)
investigator responsibilities. Process with Claude Code Opus 4.6 (1M token context).]

\subsection{Site Systems Integration}

[PROCESSING INSTRUCTION: Generate 0.75 page with Table 19 covering site systems integration.
Include EHR (FHIR R4), PACS (DICOM), LIS (HL7 v2), pharmacy, and MCP endpoints.
Reference the national MCP server infrastructure (5 servers, 23 tools).
Process with Claude Code Opus 4.6 (1M token context).]

\subsection{Patient-Facing Interface}

[PROCESSING INSTRUCTION: Generate 0.5-0.75 page on the patient interface. Cover eConsent,
scheduling, ePRO, patient education, and DHT integration. Reference the patient
instructions paper from the national platform.
Process with Claude Code Opus 4.6 (1M token context).]
